\documentclass[UTF8, 12pt, fontset=fandol, oneside]{ctexart}
\usepackage{researchnotes}

% 保留分册独立编译时的章号前缀。
\renewcommand{\thesection}{4.\arabic{section}}

\title{第四章\quad 导数与微分}
\author{}
\date{}

\begin{document}

\maketitle

人们所说的 ``微积分'' 实际上包括 ``微分学'' 和 ``积分学'' 两部分. 这一章我们就来讲述一元函数的微分学。

\section{导数}

\subsection{导数概念的引入}

导数是微分学的核心, 历史上这一概念是由对如下两个问题的研究而产生的:
\begin{enumerate}
  \item 求变速直线运动的瞬时速度;
  \item 求曲线上一点处的切线。
\end{enumerate}
这两个问题的解决最终都归结为求一个函数的变化率的问题. 牛顿 (Newton) 从第一个问题出发, 莱布尼茨 (Leibniz) 从第二个问题出发, 分别给出了导数的概念。

\noindent\textbf{1. 求变速直线运动的瞬时速度}

通常人们所说的物体的运动速度, 一般是指物体在某段时间内的平均速度. 事实上, 任何物体的运动都很难保持匀速状态, 其速度每时每刻都在变化着. 像火箭升空, 火车行驶, 都是从静止慢慢加速而达到某种速度的. 此外, 即使以一定速度运动的物体, 也会根据运动过程中遇到的各种情况, 随时调整它的运行速度. 因此, 一般来说平均速度只是对运动物体快慢程度的粗略描述, 并不能真正反映物体每一时刻运动的快慢程度. 随着科学技术的发展, 仅仅知道运动物体的平均速度是不够的, 而是需要精确地给出其每一时刻的速度 --- 瞬时速度. 那么, 如何来计算瞬时速度呢? 一般来说, 人们易于测量的是沿直线运动的物体离开初始位置的距离, 因此一个自然的想法就是利用物体所走过的路程来求瞬时速度。

设物体沿直线朝一个方向作非匀速运动, 其运动规律由函数 $s=s(t)\ (t\in[0,T])$ 给出, 其中 $t$ 是时间, $s$ 是物体在时间 $t$ 内所走过的路程. 现在要求它在时刻 $t_0\in(0,T)$ 的瞬时速度 $v(t_0)$。

显然, 用 $[0,T]$ 上的平均速度来代替 $v(t_0)$, 一般来说误差应很大. 一个朴素的想法是, $v(t_0)$ 的大小应该只与 $t_0$ 附近物体移动的路程的改变量有关. 为此, 给 $t_0$ 一个微小的改变量 $\Delta t\ne0$, 物体在 $\Delta t$ 这段时间内所经过的路程为
\[
  \Delta s=s(t_0+\Delta t)-s(t_0).
\]
物体在 $\Delta t$ 这段时间内的平均速度为
\[
  v=\frac{s(t_0+\Delta t)-s(t_0)}{\Delta t}.
\]
容易看出, $\Delta t$ 越接近于 $0$, 该平均速度就应该越接近瞬时速度 $v(t_0)$. 于是, 在时刻 $t_0\in(0,T)$ 的瞬时速度 $v(t_0)$ 就应该是当 $\Delta t\to0$ 时, 平均速度 $v$ 的极限, 即
\[
  v(t_0)
  =
  \lim_{\Delta t\to0}\frac{\Delta s}{\Delta t}
  =
  \lim_{\Delta t\to0}
  \frac{s(t_0+\Delta t)-s(t_0)}{\Delta t}.
\]
上式既给出了瞬时速度的定义, 也给出了计算瞬时速度的方法。

\noindent\textbf{2. 求曲线上一点处切线}

莱布尼茨在研究如何求曲线在一点的切线这一几何问题时, 同样遇到了类似的极限问题. 现假定有一条平面曲线, 它是函数 $y=f(x)\ (x\in(a,b))$ 的图像, 并给定 $x_0\in(a,b)$. 那么, 应该怎样来求此曲线在点 $P(x_0,f(x_0))$ 的切线呢 (参见图 4.1.1)?

由于所求的切线经过点 $(x_0,f(x_0))$, 因此我们只要将其斜率确定即可. 如同计算瞬时速度一样, 一个显然的事实是过该点的切线的斜率应该只与 $x_0$ 附近的 $x$ 所对应的函数值有关. 为此, 我们在 $x_0$ 附近取一点 $x$, 并且观察过曲线上两点 $P(x_0,f(x_0))$ 和 $Q(x,f(x))$ 的割线. 显然, 该割线的斜率为
\[
  \tilde{k}(x)=\frac{f(x)-f(x_0)}{x-x_0}.
\]
当 $x$ 在 $x_0$ 附近变化时, 点 $Q$ 在曲线上变动, 割线的斜率也在变化. 当 $x$ 趋于 $x_0$ 时, 点 $Q$ 沿曲线趋近于点 $P$, 割线 $PQ$ 就应该趋近于过点 $P$ 的切线 $PT$。

因此, 如果令 $\Delta x=x-x_0,\Delta y=f(x)-f(x_0)$, 当 $\Delta x$ 趋于 $0$ 时, 我们自然定义割线 $PQ$ 的极限位置为该曲线过点 $P$ 的切线 $PT$. 这样, 割线 $PQ$ 的斜率 $\tilde{k}(x)=\dfrac{\Delta y}{\Delta x}$ 的极限就是所求切线 $PT$ 的斜率, 即
\[
  k=
  \lim_{\Delta x\to0}\tilde{k}(x)
  =
  \lim_{\Delta x\to0}\frac{\Delta y}{\Delta x}
  =
  \lim_{x\to x_0}\frac{f(x)-f(x_0)}{x-x_0}.
\]
上面的讨论既给出了切线的定义, 也给出了计算切线斜率的方法. 上述两个问题的实际意义完全不同: 一个是运动学中沿直线运动的物体的瞬时速度, 一个是几何学中曲线在一点处的切线. 但从数学上来看, 它们是完全一样的, 即都是归结为求函数的改变量 $\Delta y$ 与自变量的改变量 $\Delta x$ 之比的极限。

\subsection{导数的定义}

除了上一小节的例子外, 还有许多实际问题也要遇到类似的求函数的改变量与自变量的改变量之比的极限问题. 例如, 由质量函数求密度, 由人口增长函数求增长率等. 略去各种背景, 便有以下纯数学的定义:

\noindent\textbf{定义 4.1.1}\quad 设函数 $y=f(x)$ 在 $U(x_0,\delta_0)$ 内有定义. 若极限
\[
  \lim_{x\to x_0}\frac{f(x)-f(x_0)}{x-x_0}
\]
存在, 则称 $f(x)$ 在点 $x_0$ 处可导, 并且称此极限值为 $f(x)$ 在 $x_0$ 处的\textbf{导数}, 记为 $f'(x_0)$ 或者 $\dfrac{df(x_0)}{dx}$, 有时也记为 $\left.\dfrac{df}{dx}\right|_{x=x_0}$ 或者 $\left.y'\right|_{x=x_0}$。

由定义即知, 函数在一点的导数是它在该点附近各点平均变化率的极限, 是函数变化率的一种定量刻画. 从几何上来看, 导数 $f'(x_0)$ 就是曲线 $y=f(x)$ 在点 $P(x_0,f(x_0))$ 的切线的斜率, 即切线与 $x$ 轴正向夹角的正切值 (参见图 4.1.1)。

\noindent\textbf{例 4.1.1}\quad 求常数函数 $f(x)=C$ 在点 $x_0\in(-\infty,+\infty)$ 的导数。

\noindent\textbf{解}\quad 对 $\forall x_0\in(-\infty,+\infty)$, 由于
\[
  \lim_{x\to x_0}\frac{f(x)-f(x_0)}{x-x_0}
  =
  \lim_{x\to x_0}\frac{C-C}{x-x_0}=0,
\]
所以 $f'(x_0)=0$。

\noindent\textbf{例 4.1.2}\quad 求 $f(x)=\sqrt{x}$ 在 $x=x_0>0$ 的导数。

\noindent\textbf{解}\quad 由于
\[
  \lim_{x\to x_0}\frac{\sqrt{x}-\sqrt{x_0}}{x-x_0}
  =
  \lim_{x\to x_0}\frac1{\sqrt{x}+\sqrt{x_0}}
  =
  \frac1{2\sqrt{x_0}},
\]
所以
\[
  f'(x_0)=\frac1{2\sqrt{x_0}}.
\]

若 $f(x)$ 在 $(a,b)$ 内每点都有导数, 则它的导数 $f'(x)$ 也是 $(a,b)$ 内的一个函数, 称为 $f(x)$ 的\textbf{导函数}, 简称导数. 此时, 如果我们要求 $f(x)$ 的导函数, 则我们先要设 $x_0$ 是任意固定的一点, 然后按定义求出极限, 再将 $x_0$ 换成 $x$. 为了避免这一过程, 我们可以引入 $\Delta x=x-x_0$, 从而 $x=x_0+\Delta x$, $f(x)=f(x_0+\Delta x)$. 因此, 我们有
\[
  f'(x)=
  \lim_{\Delta x\to0}\frac{f(x+\Delta x)-f(x)}{\Delta x}.
\]
这里 $\Delta x$ 是自变量的一个改变量, 而 $f(x+\Delta x)-f(x)$ 是函数的一个相应改变量, 记之为 $\Delta y$, 便有
\[
  f'(x)=\lim_{\Delta x\to0}\frac{\Delta y}{\Delta x}.
\]
初学者应注意的是, 这里的 $\Delta x$ 是可正可负的, 而不是恒大于 $0$ 的。

在推导基本初等函数的导数时, 我们需要如下的三个重要极限:
\[
  \text{(1) }\lim_{x\to0}\frac{\log_a(1+x)}{x}
  =
  \log_a e
  =
  \frac1{\ln a}\quad(a>0\text{且}a\ne1);
\]
\[
  \text{(2) }\lim_{x\to0}\frac{a^x-1}{x}
  =
  \ln a\quad(a>0);
\]
\[
  \text{(3) }\lim_{x\to0}\frac{(1+x)^\mu-1}{x}
  =
  \mu.
\]
在上一章中我们已经给出这些结果的详细推导过程。

\noindent\textbf{例 4.1.3}\quad 求 $f(x)=x^\alpha$ 的导数 (定义域与 $\alpha$ 有关, 对任何 $\alpha$, $x>0$ 总有定义)。

\noindent\textbf{解}\quad 设 $x\ne0$, 则有
\[
  \frac{f(x+\Delta x)-f(x)}{\Delta x}
  =
  \frac{(x+\Delta x)^\alpha-x^\alpha}{\Delta x}
  =
  x^{\alpha-1}
  \frac{\left(1+\dfrac{\Delta x}{x}\right)^\alpha-1}
       {\dfrac{\Delta x}{x}}.
\]
对固定的 $x\ne0$, 由于当 $\Delta x\to0$ 时, $\dfrac{\Delta x}{x}\to0$, 从而推出
\[
  f'(x)=\alpha x^{\alpha-1}.
\]

特别地, 我们有
\[
  (x^{-1})'=-\frac1{x^2}\quad(x\ne0),
\]
\[
  (\sqrt{x})'=\frac1{2\sqrt{x}}\quad(x>0).
\]

\noindent\textbf{例 4.1.4}\quad 求 $f(x)=a^x\ (a>0,-\infty<x<+\infty)$ 的导数。

\noindent\textbf{解}\quad 由于 $\forall x\in(-\infty,+\infty)$, 有
\[
  \frac{a^{x+\Delta x}-a^x}{\Delta x}
  =
  a^x\frac{a^{\Delta x}-1}{\Delta x}
  \to a^x\ln a\quad(\Delta x\to0),
\]
所以我们有 $f'(x)=a^x\ln a$. 特别地, 我们有 $(e^x)'=e^x$。

\noindent\textbf{例 4.1.5}\quad 求 $f(x)=\log_a x\ (0<a\ne1,0<x<+\infty)$ 的导数。

\noindent\textbf{解}\quad 由于 $\forall x\in(0,+\infty)$, 有
\[
  \frac{\log_a(x+\Delta x)-\log_a x}{\Delta x}
  =
  \frac1x
  \frac{\log_a\left(1+\dfrac{\Delta x}{x}\right)}
       {\dfrac{\Delta x}{x}}
  \to
  \frac{\log_a e}{x}\quad(\Delta x\to0),
\]
所以我们有
\[
  f'(x)=\frac{\log_a e}{x}=\frac1{x\ln a}.
\]
特别地, 我们有 $(\ln x)'=\dfrac1x$。

\noindent\textbf{例 4.1.6}\quad 求 $f(x)=\sin x\ (-\infty<x<+\infty)$ 的导数。

\noindent\textbf{解}\quad 由于 $\forall x\in(-\infty,+\infty)$, 有
\[
  \frac{\sin(x+\Delta x)-\sin x}{\Delta x}
  =
  \frac{2\cos\left(x+\dfrac{\Delta x}{2}\right)\sin\dfrac{\Delta x}{2}}
       {\Delta x}
  \to
  \cos x\quad(\Delta x\to0),
\]
所以我们有 $f'(x)=\cos x$。

\noindent\textbf{例 4.1.7}\quad 求 $f(x)=\cos x\ (-\infty<x<+\infty)$ 的导数。

\noindent\textbf{解}\quad 由于 $\forall x\in(-\infty,+\infty)$, 有
\[
  \frac{\cos(x+\Delta x)-\cos x}{\Delta x}
  =
  \frac{-2\sin\left(x+\dfrac{\Delta x}{2}\right)\sin\dfrac{\Delta x}{2}}
       {\Delta x}
  \to
  -\sin x\quad(\Delta x\to0),
\]
所以我们有 $f'(x)=-\sin x$。

\subsection{单侧导数}

类似于单侧极限, 我们有以下单侧导数的概念。

\noindent\textbf{定义 4.1.2}\quad 设函数 $y=f(x)$ 在 $U^+(x_0,\delta)$ 内有定义. 若极限
\[
  \lim_{\Delta x\to0+0}
  \frac{f(x_0+\Delta x)-f(x_0)}{\Delta x}
\]
存在, 则称 $f(x)$ 在 $x_0$ 处右可导, 并且称此极限值为 $f(x)$ 在 $x_0$ 处的\textbf{右导数}, 记为 $f'_+(x_0)$. 同理可定义函数 $f(x)$ 在 $x_0$ 处的\textbf{左可导}和\textbf{左导数}, 即
\[
  f'_-(x_0)=
  \lim_{\Delta x\to0-0}
  \frac{f(x_0+\Delta x)-f(x_0)}{\Delta x}.
\]

\noindent\textbf{注}\quad 从导数的定义我们可以看出, $f'(x_0)$ 存在的充分必要条件是 $f'_+(x_0)$ 与 $f'_-(x_0)$ 都存在且相等。

当 $f(x)$ 在 $x_0$ 处可导时, 有
\[
  f'(x_0)=
  \lim_{\Delta x\to0}
  \frac{f(x_0+\Delta x)-f(x_0)}{\Delta x}.
\]
因此, $f'(x_0)$ 存在的一个必要条件是, 当 $\Delta x\to0$ 时, $\Delta y=f(x+\Delta x)-f(x)$ 必须趋于零. 而这正是 $f(x)$ 在 $x_0$ 处连续的定义, 从而我们有

\noindent\textbf{定理 4.1.1}\quad 若函数 $y=f(x)$ 在 $x_0$ 处可导, 则它在 $x_0$ 处连续。

这样, 一个自然的问题就是, 是否连续函数一定可导呢? 其答案是否定的. 例如, 函数 $y=f(x)=|x|$ 在 $x=0$ 处是连续的, 但 $f(x)$ 在 $x=0$ 处不可导. 这是由于当 $\Delta x\to0$ 时,
\[
  \frac{f(0+\Delta x)-f(0)}{\Delta x}
  =
  \frac{|\Delta x|}{\Delta x}
  =
  \begin{cases}
    1, & \Delta x>0,\\
    -1, & \Delta x<0
  \end{cases}
\]
的极限不存在 (参见图 4.1.2)。

\noindent\textbf{注}\quad 从单侧导数的定义可以看出, $f(x)$ 在 $x_0$ 处左可导, 则 $f(x)$ 在 $x_0$ 处左连续; 右可导则右连续. 于是 $f(x)$ 在 $x_0$ 左、右可导 (不要求左导数与右导数相等) 就可以推出 $f(x)$ 在 $x_0$ 连续。

\noindent\textbf{例 4.1.8}\quad 令函数
\[
  S(x)=
  \begin{cases}
    \sin\dfrac1x, & x\ne0,\\
    0, & x=0,
  \end{cases}
\]
则有
\begin{enumerate}
  \item $S(x)$ 在 $x=0$ 处间断, 是第二类间断点;
  \item $xS(x)$ 在 $x=0$ 处连续, 但不可导 (其实, 其单侧导数也不存在);
  \item $x^2S(x)$ 在 $x=0$ 处连续而且可导, 其导数为 $0$。
\end{enumerate}
此外, 按照导数的定义, 曲线 $y=x^2\sin\dfrac1x$ 在 $(0,0)$ 处的切线方程为 $y=0$, 此时, 该切线在 $x=0$ 附近总是与曲线有交点, 这与中学里曲线的切线定义是有很大区别的。

\section{求导数的方法}

按照导数定义, 导数表示为当 $\Delta x\to0$ 时, $\dfrac{\Delta y}{\Delta x}$ 的极限. 但直接用定义对各种函数求导数, 并不总是一件容易的事, 有时甚至很难办到. 因此, 这一节我们来介绍各种行之有效的求导方法. 利用这些求导方法就可以轻而易举地求出许多函数的导数。

\subsection{函数四则运算的导数}

\noindent\textbf{定理 4.2.1}\quad 设函数 $f(x)$ 和 $g(x)$ 都在点 $x$ 处可导, 则下列各式在点 $x$ 处成立:
\[
  \text{(1) }\quad
  (f(x)\pm g(x))'=f'(x)\pm g'(x);
\]
\[
  \text{(2) }\quad
  (f(x)g(x))'=f'(x)g(x)+f(x)g'(x);
\]
\[
  \text{(3) }\quad
  \left(\frac{f(x)}{g(x)}\right)'
  =
  \frac{f'(x)g(x)-f(x)g'(x)}{g^2(x)}
  \quad(g(x)\ne0).
\]

\noindent\textbf{证明}\quad (1) 令 $y=f(x)\pm g(x)$, 则有
\[
  (f(x)\pm g(x))'
  =
  \lim_{\Delta x\to0}\frac{\Delta y}{\Delta x}
\]
\[
  =
  \lim_{\Delta x\to0}
  \frac{[f(x+\Delta x)\pm g(x+\Delta x)]-[f(x)\pm g(x)]}
       {\Delta x}
\]
\[
  =
  \lim_{\Delta x\to0}\frac{f(x+\Delta x)-f(x)}{\Delta x}
  \pm
  \lim_{\Delta x\to0}\frac{g(x+\Delta x)-g(x)}{\Delta x}
  =
  f'(x)\pm g'(x).
\]

(2) 令 $y=f(x)g(x)$, 则有
\[
  (f(x)g(x))'
  =
  \lim_{\Delta x\to0}\frac{\Delta y}{\Delta x}
\]
\[
  =
  \lim_{\Delta x\to0}
  \frac{f(x+\Delta x)g(x+\Delta x)-f(x)g(x)}
       {\Delta x}
\]
\[
  =
  \lim_{\Delta x\to0}
  \frac{f(x+\Delta x)-f(x)}{\Delta x}g(x+\Delta x)
  +
  \lim_{\Delta x\to0}
  f(x)\frac{g(x+\Delta x)-g(x)}{\Delta x}
\]
\[
  =
  f'(x)g(x)+f(x)g'(x).
\]

(3) 令 $y=\dfrac1{g(x)}$, 则有
\[
  y'=
  \lim_{\Delta x\to0}\frac{\Delta y}{\Delta x}
\]
\[
  =
  \lim_{\Delta x\to0}
  \frac1{\Delta x}
  \left[
    \frac1{g(x+\Delta x)}-\frac1{g(x)}
  \right]
\]
\[
  =
  -
  \lim_{\Delta x\to0}
  \frac{g(x+\Delta x)-g(x)}{\Delta x}
  \frac1{g(x+\Delta x)g(x)}
  =
  -\frac{g'(x)}{g^2(x)}.
\]
再由 $\dfrac{f(x)}{g(x)}=f(x)\dfrac1{g(x)}$ 及 (2), 我们有
\[
  \left(\frac{f(x)}{g(x)}\right)'
  =
  \frac{f'(x)}{g(x)}
  +
  f(x)\left(\frac1{g(x)}\right)'
  =
  \frac{f'(x)g(x)-f(x)g'(x)}{g^2(x)}.
\]

\noindent\textbf{推论}\quad 设 $f_i(x)\ (i=1,2,\cdots,n)$ 都在点 $x$ 处可导, 则下列各式在点 $x$ 处成立:
\[
  \text{(1) }\quad
  \left(\sum_{i=1}^{n}c_if_i(x)\right)'
  =
  \sum_{i=1}^{n}c_if_i'(x),
  \quad\text{其中 }c_i\text{ 是常数};
\]
\[
  \text{(2) }\quad
  \left(\prod_{i=1}^{n}f_i(x)\right)'
  =
  \sum_{k=1}^{n}
  \left(
    f_k'(x)\prod_{\substack{i=1\\i\ne k}}^{n}f_i(x)
  \right).
\]

\noindent\textbf{例 4.2.1}\quad 设 $f(x)=\tan x$, 求 $f'(x)$。

\noindent\textbf{解}\quad
\[
  f'(x)=(\tan x)'=\left(\frac{\sin x}{\cos x}\right)'
\]
\[
  =
  \frac{(\sin x)'\cos x-\sin x(\cos x)'}{\cos^2 x}
  =
  \frac{\cos^2 x+\sin^2 x}{\cos^2 x}
  =
  \frac1{\cos^2 x}
  =
  \operatorname{sec}^2x.
\]

\noindent\textbf{例 4.2.2}\quad 设 $f(x)=\cot x$, 求 $f'(x)$。

\noindent\textbf{解}\quad
\[
  (\cot x)'=\left(\frac1{\tan x}\right)'
  =
  -\frac{(\tan x)'}{\tan^2 x}
\]
\[
  =
  -\frac1{\cos^2 x}\frac{\cos^2 x}{\sin^2 x}
  =
  -\frac1{\sin^2 x}
  =
  -\operatorname{csc}^2x.
\]

\subsection{反函数的求导法则}

\noindent\textbf{定理 4.2.2}\quad 设函数 $y=f(x)$ 在 $(a,b)$ 内严格单调, 并且令
\[
  \alpha=\min\{f(a+0),f(b-0)\},
  \quad
  \beta=\max\{f(a+0),f(b-0)\}.
\]
如果 $f(x)$ 在 $(a,b)$ 内可导且导数 $f'(x)\ne0$, 则它的反函数 $x=f^{-1}(y)$ 在 $(\alpha,\beta)$ 内可导, 而且有
\[
  \frac{df^{-1}(y)}{dy}=\frac1{f'(x)}.
\]

\noindent\textbf{证明}\quad 由于 $y=f(x)$ 在 $(a,b)$ 内可导, 从而它在 $(a,b)$ 上连续. 再由其在 $(a,b)$ 内严格单调知, 它的反函数在 $(\alpha,\beta)$ 中连续且严格单调. 任取 $y_0\in(\alpha,\beta)$, 则存在 $x_0\in(a,b)$, 使得 $y_0=f(x_0)$, 而且
\[
  \left.(f^{-1}(y))'\right|_{y=y_0}
  =
  \lim_{y\to y_0}\frac{f^{-1}(y)-f^{-1}(y_0)}{y-y_0}
\]
\[
  =
  \lim_{y\to y_0}\frac{x-x_0}{y-y_0}
  =
  \lim_{x\to x_0}
  \frac1{\dfrac{f(x)-f(x_0)}{x-x_0}}
  =
  \frac1{f'(x_0)}.
\]

\noindent\textbf{例 4.2.3}\quad 求 $y=\arctan x$ 的导数。

\noindent\textbf{解}\quad 由 $y=\arctan x$, 得 $x=\tan y$. 因此, 有
\[
  y'=(\arctan x)'=\frac1{(\tan y)'}
  =\frac1{\sec^2 y}
  =\frac1{1+\tan^2 y}
  =\frac1{1+x^2}.
\]
类似地, 有
\[
  (\operatorname{arccot}x)'=-\frac1{1+x^2}.
\]

\noindent\textbf{例 4.2.4}\quad 求 $y=\arcsin x$ 的导数。

\noindent\textbf{解}\quad 由 $x=\sin y$, 得
\[
  y'=(\arcsin x)'=\frac1{(\sin y)'}
  =\frac1{\cos y}
  =\frac1{\sqrt{1-\sin^2y}}
  =\frac1{\sqrt{1-x^2}}.
\]
类似地, 有
\[
  (\arccos x)'=-\frac1{\sqrt{1-x^2}}.
\]

在利用反函数求导法时, 读者需注意两点: 一是要注意每次求导是关于什么变量进行的; 二是如何将求导数后关于 $y$ 的函数转化为 $x$ 的函数。

通过一系列例题, 我们已经求出了所有基本初等函数的导数. 现将这些结果总结如下, 以便以后查阅:
\[
\begin{array}{ll}
  \text{(1) }(C)'=0; &
  \text{(2) }(x^\alpha)'=\alpha x^{\alpha-1};\\[4pt]
  \text{(3) }(\sin x)'=\cos x; &
  \text{(4) }(\cos x)'=-\sin x;\\[4pt]
  \text{(5) }(\tan x)'=\dfrac1{\cos^2x}; &
  \text{(6) }(\cot x)'=-\dfrac1{\sin^2x};\\[8pt]
  \text{(7) }(\arcsin x)'=\dfrac1{\sqrt{1-x^2}}; &
  \text{(8) }(\arccos x)'=-\dfrac1{\sqrt{1-x^2}};\\[8pt]
  \text{(9) }(\arctan x)'=\dfrac1{1+x^2}; &
  \text{(10) }(\operatorname{arccot}x)'=-\dfrac1{1+x^2};\\[8pt]
  \text{(11) }(e^x)'=e^x; &
  \text{(12) }(a^x)'=a^x\ln a\ (a>0);\\[4pt]
  \text{(13) }(\ln x)'=\dfrac1x; &
  \text{(14) }(\log_a x)'=\dfrac1{x\ln a}\quad\text{若 }(a>0\text{ 且 }a\ne1).
\end{array}
\]

\subsection{复合函数的求导法则}

\noindent\textbf{定理 4.2.3}\quad 设函数 $y=f(u)$ 在 $U(u_0,\delta_0)$ 内有定义, 函数 $u=g(x)$ 在 $U(x_0,\eta_0)$ 内有定义, 且 $u_0=g(x_0)$. 若 $f'(u_0)$ 与 $g'(x_0)$ 都存在, 则复合函数 $F(x)=f(g(x))$ 在点 $x_0$ 可导, 且 $F'(x_0)=f'(g(x_0))g'(x_0)$。

\noindent\textbf{证明}\quad 定义函数
\[
  A(u)=
  \begin{cases}
    \dfrac{f(u)-f(u_0)}{u-u_0}, & u\ne u_0,\\[8pt]
    f'(u_0), & u=u_0,
  \end{cases}
\]
则 $A(u)$ 在 $u_0$ 点连续, 即有
\[
  \lim_{u\to u_0}A(u)=A(u_0)=f'(u_0).
\]
由恒等式 $f(u)-f(u_0)=A(u)(u-u_0)$, 我们有
\[
  \frac{F(x)-F(x_0)}{x-x_0}
  =
  \frac{f(g(x))-f(g(x_0))}{x-x_0}
  =
  A(g(x))\frac{g(x)-g(x_0)}{x-x_0}.
\]
在上式两边令 $x\to x_0$, 得
\[
  F'(x_0)=f'(g(x_0))g'(x_0).
\]

\noindent\textbf{注}\quad 若 $f(u)$ 的定义域包含 $u=g(x)$ 的值域, 两函数在各自的定义域上可导, 则复合函数 $F(x)=f(g(x))$ 在 $g(x)$ 的定义域上可导, 且
\[
  F'(x)=f'(g(x))g'(x).
\]
上述等式常常简记为
\[
  y_x'=y_u'u_x'
  \quad\text{或者}\quad
  \frac{dy}{dx}=\frac{dy}{du}\frac{du}{dx}.
\]
我们一般称它为链锁法则。

\noindent\textbf{例 4.2.5}\quad 分别求 $y=\sin x^2$ 和 $y=\sin^2x$ 的导数。

\noindent\textbf{解}\quad $y=\sin x^2$ 可看成 $y=\sin u$ 和 $u=x^2$ 的复合, 因此
\[
  (\sin x^2)'=2x\cos x^2;
\]
而 $y=\sin^2x$ 可看成 $y=u^2$ 和 $u=\sin x$ 的复合, 因此
\[
  (\sin^2x)'=2\sin x(\sin x)'=2\sin x\cos x=\sin2x.
\]

\noindent\textbf{例 4.2.6}\quad 求 $y=\ln|x|\ (x\ne0)$ 的导数。

\noindent\textbf{解}\quad 当 $x>0$ 时, $y=\ln x$, 从而
\[
  y'=(\ln x)'=\frac1x;
\]
当 $x<0$ 时,
\[
  y=\ln|x|=\ln(-x),
\]
于是, 它可看成由 $y=\ln u$ 与 $u=-x$ 复合而成, 从而有
\[
  (\ln|x|)'=(\ln(-x))'=\frac1{-x}(-x)'=\frac1x.
\]

\noindent\textbf{例 4.2.7}\quad 求 $y=x^\alpha\ (x>0,\alpha\text{ 为任意实数})$ 的导数。

\noindent\textbf{解}\quad 前面我们用定义已求出 $y'=\alpha x^{\alpha-1}$. 现应用复合函数求导法于 $y=e^{\alpha\ln x}$, 则有
\[
  y'=(e^{\alpha\ln x})'
  =
  e^{\alpha\ln x}(\alpha\ln x)'
  =
  \frac{\alpha e^{\alpha\ln x}}x
  =
  \alpha\frac{x^\alpha}x
  =
  \alpha x^{\alpha-1}.
\]

\noindent\textbf{例 4.2.8}\quad 证明: 当 $n\ge2$ 时, 有
\[
  \sum_{k=1}^{n}C_n^kk^2=n(n+1)2^{n-2}.
\]

\noindent\textbf{证明}\quad 在恒等式
\[
  (1+x)^n=\sum_{k=0}^{n}C_n^kx^k
\]
两边对 $x$ 求导数得
\[
  n(1+x)^{n-1}=\sum_{k=1}^{n}C_n^kkx^{k-1}.
\]
在上述恒等式的两边乘以 $x$ 并对 $x$ 求导数得
\[
  n(1+x)^{n-1}+n(n-1)x(1+x)^{n-2}
  =
  \sum_{k=1}^{n}C_n^kk^2x^{k-1}.
\]
在上述等式中令 $x=1$ 整理得
\[
  n(n+1)2^{n-2}=\sum_{k=1}^{n}C_n^kk^2.
\]

\subsection{隐函数的求导法}

前面我们所遇到的函数都是由显式给出的. 但在解决实际问题时, 常常得到的只是变量 $x$ 和 $y$ 所满足的方程, 如:
\[
  x^2+y^2=1,\quad \tan x+\tan y=xy,\quad \ln y+\sin y=x^2y,
\]
等等. 通常在一定条件下, 对指定数集 $X$ 中的每一个数 $x$, 由变量 $x$ 和 $y$ 所满足的方程可以唯一地确定一个数 $y$ 与之对应. 由这种对应法则所确定的函数 $y=y(x)$ 就称做由该方程所确定的隐函数. 例如, 方程 $x^2+y^2=1$ 就隐含了两个定义在 $[-1,1]$ 上的函数 $y=f_1(x)=\sqrt{1-x^2}$ 和 $y=f_2(x)=-\sqrt{1-x^2}$. 至于在什么条件下才能保证由所给出的二元方程唯一地确定一个连续可导的隐函数呢? 这要等到学习了多元函数的微分学才能给出解答. 在本节的后继部分, 我们总是假定给定的二元方程所确定函数是可导的, 然后来讨论该函数的求导问题。

\noindent\textbf{例 4.2.9}\quad 证明开普勒 (Kepler) 方程
\[
  y-x-\varepsilon\sin y=0
  \quad(0<\varepsilon<1)
\]
可确定一个隐函数 $y=y(x)$, 并且求 $y'(x)$。

\noindent\textbf{证明}\quad 令 $x=\varphi(y)=y-\varepsilon\sin y$, 则有
\[
  \lim_{y\to-\infty}\varphi(y)=-\infty,
  \quad
  \lim_{y\to+\infty}\varphi(y)=+\infty;
\]
而且对任意的 $y_1<y_2$, 有
\[
  \varphi(y_2)-\varphi(y_1)
  =
  y_2-y_1-\varepsilon(\sin y_2-\sin y_1)
\]
\[
  =
  y_2-y_1-\varepsilon\left(2\cos\frac{y_2+y_1}{2}\sin\frac{y_2-y_1}{2}\right)
  \ge(1-\varepsilon)(y_2-y_1)>0.
\]
因此, $\varphi(y)$ 是一严格单调递增的连续函数, 从而 $\forall x_0\in(-\infty,+\infty)$, 总存在唯一 $y_0$, 使得 $x_0=y_0-\varepsilon\sin y_0$. 这表明, 该方程确定了 $(-\infty,+\infty)$ 上的一个函数 $y=y(x)$. 以后我们还将知道, 开普勒方程确定的函数是可导的. 下面我们来求出此函数的导数。

既然 $y=y(x)$ 是由开普勒方程确定的函数, 则该函数满足恒等式
\[
  y(x)-x-\varepsilon\sin y(x)\equiv0.
\]
两边对 $x$ 求导数后得恒等式
\[
  y'(x)-1-\varepsilon\cos y(x)y'(x)\equiv0,
\]
从而有
\[
  y'(x)=\frac1{1-\varepsilon\cos y(x)}.
\]

这一例子表明, 虽然我们并没有解出函数的显式表达式, 但我们可以由它所满足的方程求出它的导数. 当然, 导数的表达式中仍含有隐函数 $y(x)$. 同时, 这一例子也告诉我们求这类函数导数的方法, 即只要将 $y$ 看成 $x$ 的函数, 在方程的两边对 $x$ 求导数, 就可得到 $y'(x)$ 满足的恒等式, 然后再从中将 $y'(x)$ 解出即可。

\noindent\textbf{例 4.2.10}\quad 给定圆的方程 $x^2+y^2=1$, 求其斜率为 $-\dfrac12$ 的切线方程。

\noindent\textbf{解}\quad 设所求的切线经过圆上一点 $(x_0,y_0)$. 由隐函数求导法, 得 $2x+2yy'=0$, 从而有 $y'(x)=-\dfrac{x}{y}$. 所以在切点 $(x_0,y_0)$ 处应有
\[
  y'(x_0)=-\frac{x_0}{y_0}=-\frac12,
\]
即 $y_0=2x_0$. 再由 $(x_0,y_0)$ 在圆上可知, $x_0=\pm\dfrac1{\sqrt5},y_0=\pm\dfrac2{\sqrt5}$. 于是, 所求的切线有两条, 其分别过 $\left(\dfrac1{\sqrt5},\dfrac2{\sqrt5}\right)$ 和 $\left(-\dfrac1{\sqrt5},-\dfrac2{\sqrt5}\right)$ 的方程为
\[
  y=\frac2{\sqrt5}-\frac12\left(x-\frac1{\sqrt5}\right)
  =
  \frac{\sqrt5}2-\frac12x
\]
和
\[
  y=-\frac2{\sqrt5}-\frac12\left(x+\frac1{\sqrt5}\right)
  =
  -\frac{\sqrt5}2-\frac12x.
\]

对某些函数求导数时, 若从显函数出发来求比较复杂或不好求, 这时可转化为隐函数来求导数, 常用的方法是对函数的两边取对数。

\noindent\textbf{例 4.2.11}\quad 设 $y=x^{e^x}\ (x>0)$, 求 $y'$。

\noindent\textbf{解}\quad 对 $y=x^{e^x}$ 两边取对数, 得 $\ln y=e^x\ln x$, 再两边对 $x$ 求导数, 得
\[
  \frac{y'}y=e^x\ln x+\frac{e^x}x
  =
  e^x\left(\frac{x\ln x+1}{x}\right),
\]
所以
\[
  y'=x^{e^x}e^x\left(\frac{x\ln x+1}{x}\right).
\]

\noindent\textbf{例 4.2.12}\quad 求
\[
  y=\left(1+\frac1{x^2}\right)^{x^2}
\]
的导数。

\noindent\textbf{解}\quad 在 $y=\left(1+\dfrac1{x^2}\right)^{x^2}$ 的两边取对数得
\[
  \ln y=x^2\ln\left(1+\frac1{x^2}\right)
  =
  x^2[\ln(x^2+1)-\ln x^2].
\]
两边对 $x$ 求导数得
\[
  \frac{y'}y
  =
  2x(\ln(x^2+1)-\ln x^2)
  +
  x^2\left(\frac{2x}{1+x^2}-\frac{2x}{x^2}\right)
\]
\[
  =
  2x\left[\ln\left(1+\frac1{x^2}\right)-\frac1{1+x^2}\right],
\]
所以
\[
  y'=2x\left(1+\frac1{x^2}\right)^{x^2}
  \left[\ln\left(1+\frac1{x^2}\right)-\frac1{1+x^2}\right].
\]

\noindent\textbf{注}\quad 以上例子中的求导方法也称为对数求导法, 它适合对一些比较复杂的函数求导。

\subsection{参数式函数的求导法}

下面我们来讨论由参数方程
\[
  \begin{cases}
    x=x(t),\\
    y=y(t),
  \end{cases}
  \quad t\in(\alpha,\beta)
\]
所确定的函数 $y=f(x)$ 的求导问题。

假定函数 $x(t)$ 和 $y(t)$ 都在 $(\alpha,\beta)$ 内可导, 且 $x'(t)\ne0$. 在这样的条件下, 下一章我们将证明 $x(t)$ 在 $(\alpha,\beta)$ 严格单调. 令
\[
  a=\min(x(\alpha+0),x(\beta-0)),
\]
\[
  b=\max(x(\alpha+0),x(\beta-0)),
\]
则 $x=x(t)$ 的反函数 $t=t(x)$ 在 $(a,b)$ 上有定义. 因此, 由该参数方程可以确定一个函数
\[
  y=f(x)=y(t(x)),\quad x\in(a,b).
\]
由复合函数及反函数的求导法则, 得
\[
  y_x'=y'(t)t'(x)=\frac{y'(t)}{x'(t)}.
\]
这表明, 由参数方程所确定的函数的导数就等于 $y$ 与 $x$ 关于参数 $t$ 的导数的商. 在以后的计算中, 我们只要记住此公式即可, 但理解公式的推导过程是十分有益的。

\noindent\textbf{例 4.2.13}\quad 求椭圆 $\dfrac{x^2}{a^2}+\dfrac{y^2}{b^2}=1\ (a,b>0)$ 上点 $(x_0,y_0)$ 处的切线方程。

\noindent\textbf{解法 1}\quad 利用隐函数求导法, 得
\[
  \frac{2x}{a^2}+\frac{2yy'}{b^2}=0,
\]
即
\[
  y_x'=-\frac{b^2x}{a^2y}.
\]
由此可知, 当 $y_0\ne0$ 时, 在 $(x_0,y_0)$ 处的切线方程为
\[
  y-y_0=-\frac{b^2x_0}{a^2y_0}(x-x_0).
\]
注意到 $(x_0,y_0)$ 在椭圆上, 最后得切线方程为
\[
  \frac{x_0x}{a^2}+\frac{y_0y}{b^2}=1.
\]
为了求曲线在点 $(a,0)$ 处切线, 在此点的附近我们可以将 $x$ 看成 $y$ 的函数, 再利用隐函数求导法, 得
\[
  \left.x_y'\right|_{y=0}
  =
  -\left.\frac{a^2y}{b^2x}\right|_{(a,0)}
  =
  0,
\]
从而在点 $(a,0)$ 处的切线方程为 $x=a$. 同理, 可得在点 $(-a,0)$ 处的切线方程为 $x=-a$。

\noindent\textbf{解法 2}\quad 椭圆的参数方程为
\[
  \begin{cases}
    x=a\cos t,\\
    y=b\sin t,
  \end{cases}
  \quad t\in[0,2\pi).
\]
由参数式函数的求导法, 得
\[
  y_x'=\frac{(b\sin t)'}{(a\cos t)'}
  =
  -\frac{b\cos t}{a\sin t}
  \quad(t\ne0,\pi),
\]
从而在 $t_0\ne0,\pi$ 处的切线方程为
\[
  y-b\sin t_0=
  \frac{-b\cos t_0}{a\sin t_0}(x-a\cos t_0).
\]
令 $x_0=a\cos t_0,y_0=b\sin t_0$, 整理即得
\[
  \frac{x_0x}{a^2}+\frac{y_0y}{b^2}=1.
\]
完全类似于解法 1, 可求出在 $(\pm a,0)$ 处的切线方程。

\subsection{极坐标式函数的求导法}

设曲线由极坐标形式
\[
  r=r(\theta),\quad \theta\in(\alpha,\beta)
\]
给出, 其中 $r$ 是极距, $\theta$ 表示极角. 现对曲线上一点 $(\theta,r(\theta))$ 来求切线的斜率。

由极坐标方程即可得曲线的参数方程为
\[
  \begin{cases}
    x=r(\theta)\cos\theta,\\
    y=r(\theta)\sin\theta.
  \end{cases}
\]
这样, 假定 $r$ 作为 $\theta$ 的函数其导数 $r'(\theta)$ 存在, 而且在所考虑的极角 $\theta$ 附近有
\[
  x'(\theta)=r'(\theta)\cos\theta-r(\theta)\sin\theta\ne0,
\]
则由参数式函数的求导法知, 极坐标表示下的曲线在点 $(\theta,r(\theta))$ 处的切线的斜率为
\[
  \frac{dy}{dx}
  =
  \frac{r'(\theta)\sin\theta+r(\theta)\cos\theta}
       {r'(\theta)\cos\theta-r(\theta)\sin\theta}
  =
  \frac{\tan\theta+\dfrac{r(\theta)}{r'(\theta)}}
       {1-\tan\theta\dfrac{r(\theta)}{r'(\theta)}}.
\]
上式给出了极坐标式函数的求导公式. 由导数几何意义知, $\dfrac{dy}{dx}$ 是切线与 $x$ 轴正向的夹角 $\alpha$ 的正切, 即
\[
  \tan\alpha=
  \frac{\tan\theta+\dfrac{r(\theta)}{r'(\theta)}}
       {1-\tan\theta\dfrac{r(\theta)}{r'(\theta)}}.
\]
由此可解出
\[
  \frac{r(\theta)}{r'(\theta)}
  =
  \frac{\tan\alpha-\tan\theta}{1+\tan\alpha\tan\theta}
  =
  \tan(\alpha-\theta).
\]
它具有鲜明的几何意义. 在极坐标系中同时建立直角坐标系, 记向径沿逆时针方向转到切线位置的角度为 $\beta$（通常称做向径与切线的夹角）, 则有
\[
  \tan\beta=\tan(\alpha-\theta)=\frac{r(\theta)}{r'(\theta)}.
\]
这是因为这二个角之间有关系式: $\beta=k\pi+(\alpha-\theta)$, 这里 $k$ 是非负整数, 且满足 $0\le k\pi+(\alpha-\theta)\le\pi$. 例如, $k=0$ 的情形如图 4.2.1 所示。

\noindent\textbf{例 4.2.14}\quad 求证: 双纽线 $r^2=a^2\cos2\theta\ (0<\theta<\pi/4)$ 的向径与切线的夹角等于极角的两倍加 $\pi/2$。

\noindent\textbf{证明}\quad 设向径与切线的夹角为 $\beta$, 则
\[
  \tan\beta
  =
  \frac{r(\theta)}{r'(\theta)}
  =
  \frac{r^2(\theta)}{r(\theta)r'(\theta)}
  =
  -\frac{\cos2\theta}{\sin2\theta}
  =
  -\cot2\theta
  =
  \tan\left(2\theta+\frac{\pi}{2}\right).
\]
注意到 $0<\theta<\dfrac{\pi}{4}$, 便有 $\beta=2\theta+\dfrac{\pi}{2}$。

最后指出, 显式表示、隐式表示、参数表示和极坐标表示是表达函数的四种重要形式, 各有不同的适用场合, 有很多函数只能用其中的某一种来表达. 即使有些函数能够用它们中的任何一种来表达, 但运算的难易程度也可能是大相径庭的, 因此必须全面地掌握相应的求导数公式, 根据实际情况选择使用。

\section{微分}

\subsection{微分的定义}

与函数在一点的可导性紧密联系着的一个概念是可微性, 本节就来讨论函数的可微性. 我们先来看一个具体的例子. 若用 $S$ 表示半径为 $r$ 的圆的面积, 则 $S$ 可以表示为 $r$ 的函数: $S=\pi r^2$. 如果给半径 $r$ 一个增量 $\Delta r$, 则面积 $S$ 相应地有增量
\[
  \Delta S
  =
  \pi(r+\Delta r)^2-\pi r^2
  =
  2\pi r\Delta r+\pi(\Delta r)^2.
\]
它是半径为 $r$ 与 $r+\Delta r$ 的两个同心圆之间的圆环的面积. 从上式可以看出, $\Delta S$ 由两部分组成: 第一部分 $2\pi r\Delta r$ 是 $\Delta r$ 的线性函数, 而第二部分 $\pi(\Delta r)^2$ 是较 $\Delta r$ 高阶的无穷小量, 即 $\pi(\Delta r)^2=o(|\Delta r|)\ (\Delta r\to0)$. 由此可知, 当给半径 $r$ 一个微小的增量 $\Delta r$ 时, 所引起的圆面积的增量 $\Delta S$ 就可以近似地用第一部分 $2\pi r\Delta r$ 来代替, 由此所产生的相对误差为
\[
  \left|\frac{\pi(\Delta r)^2}{2\pi r\Delta r}\right|
  =
  \frac{|\Delta r|}{2r}.
\]
显然, $\Delta r$ 越小, 所产生的相对误差也越小。

对于一般的情形, 我们引入如下的定义:

\noindent\textbf{定义 4.3.1}\quad 设函数 $y=f(x)$ 在 $U(x_0,\delta_0)$ 内有定义. 如果存在常数 $A$, 使得
\[
  \Delta y=f(x_0+\Delta x)-f(x_0)=A\Delta x+o(\Delta x)
  \quad(\Delta x\to0),
\]
则称 $f(x)$ 在点 $x_0$ 处可微, 并称 $A\Delta x$ 为 $f(x)$ 在 $x_0$ 处的微分, 记做
\[
  dy=A\Delta x\quad\text{或者}\quad df(x_0)=A\Delta x.
\]

从定义看到, 一个函数的微分具有两个特点:
\begin{enumerate}
  \item 它是自变量的增量 $\Delta x$ 的线性函数。
  \item 它与函数的增量 $\Delta y$ 之差是较 $\Delta x$ 高阶的无穷小量 $(\Delta x\to0)$。
\end{enumerate}
根据这一特点, 我们就可以用 $dy$ 来近似地表达 $\Delta y$, 所产生的相对误差是一个无穷小量 $(\Delta x\to0)$, 而且 $|\Delta x|$ 越小, 近似程度就越好. 因此, 我们称 $dy=A\Delta x$ 为 $\Delta y$ 的\textbf{线性主部}. 此外, 由于 $dy$ 是 $\Delta x$ 的线性函数, 所以它特别容易计算。

微分与导数有着密切的关系, 表现为如下定理:

\noindent\textbf{定理 4.3.1}\quad 函数 $f(x)$ 在点 $x_0$ 处可微的充分必要条件是 $f(x)$ 在 $x_0$ 处可导。

\noindent\textbf{证明}\quad 必要性\quad 设 $f(x)$ 在 $x_0$ 处可微, 即
\[
  \Delta y=f(x_0+\Delta x)-f(x_0)=A\Delta x+o(\Delta x)
  \quad(\Delta x\to0)
\]
成立, 则有
\[
  \frac{\Delta y}{\Delta x}
  =
  A+o(1)
  \quad(\Delta x\to0),
\]
从而
\[
  f'(x_0)
  =
  \lim_{\Delta x\to0}\frac{\Delta y}{\Delta x}
  =
  A.
\]
这表明, $f(x)$ 在 $x_0$ 处可导, 且 $f'(x_0)=A$。

充分性\quad 设 $f(x)$ 在 $x_0$ 处可导, 即
\[
  \lim_{\Delta x\to0}
  \frac{f(x_0+\Delta x)-f(x_0)}{\Delta x}
  =
  f'(x_0)
\]
成立, 于是有
\[
  \frac{f(x_0+\Delta x)-f(x_0)}{\Delta x}
  =
  f'(x_0)+o(1)
  \quad(\Delta x\to0),
\]
即
\[
  \Delta y=f(x_0+\Delta x)-f(x_0)
  =
  f'(x_0)\Delta x+o(\Delta x)
  \quad(\Delta x\to0).
\]
这表明, $f(x)$ 在点 $x_0$ 处可微, 且 $df=f'(x_0)\Delta x$。

\noindent\textbf{注}\quad 由于这一定理的缘故, 对于一元函数而言, 人们常把 ``可导'' 和 ``可微'' 等同起来使用, 求导方法又称为微分法. 此外, 从定理的证明过程知, 微分定义式中的系数 $A$ 唯一地确定, 而且必有 $A=f'(x_0)$ 成立。

现考虑一个具体函数 $y=x$ 的微分. 由定义, 得 $dy=dx=1\Delta x$, 即 $dx=\Delta x$. 因此, 我们就可在微分记号中将 $\Delta x$ 改为 $dx$, 即若 $y=f(x)$ 在 $x_0$ 处可微, 则 $dy=f'(x_0)dx$. 因此, 导数的记号
\[
  \left.\frac{dy}{dx}\right|_{x=x_0}=f'(x_0)
\]
就可以看成 $dy$ 与 $dx$ 的比值了, 即导数乃是微分之商, 故又称微商。

下面我们来看微分的几何意义. 设函数 $y=f(x)$ 的图像如图 4.3.1 所示, 令 $N=(x_0+\Delta x,f(x_0))$, $P=(x_0+\Delta x,f(x_0+\Delta x))$, $K=(x_0+\Delta x,f(x_0)+f'(x_0)\Delta x)$, 并且过 $M=(x_0,f(x_0))$ 作曲线的切线, 则有
\[
\begin{gathered}
  dx=\Delta x=MN\text{ 为自变量的增量;}\\
  \Delta y=PN\text{ 为函数的增量;}\\
  dy=f'(x_0)dx=KN\text{ 为切线的增量;}\\
  \Delta y-dy=o(\Delta x)=PK\text{ 为函数的增量与切线的增量之差.}
\end{gathered}
\]
这表明, 微分就是在点 $x_0$ 的附近以切线的增量拟合曲线的增量. 因此, 所谓微分的基本原理就是: 一个可微函数, 在每一点的充分小局部, 均可近似视为增量按比例变化的线性函数。

若函数 $y=f(x)$ 在区间 $(a,b)$ 内每一点 $x$ 处可微 (即 $dy=f'(x)dx$), 则称函数 $f(x)$ 在 $(a,b)$ 内可微。

利用微分可作误差估计和近似计算, 下面举两个简单的例子。

\noindent\textbf{例 4.3.1}\quad 求 $\sqrt[4]{80}$ 的近似值。

\noindent\textbf{解}\quad 取 $f(x)=\sqrt[4]{x}$, 令 $x_0=81,\Delta x=-1$, 则有
\[
  f'(x)=\frac{1}{4\sqrt[4]{x^3}},
  \quad
  f(x_0)=3,
  \quad
  f'(x_0)=\frac{1}{108}.
\]
于是
\[
  \sqrt[4]{80}\approx 3+\frac{1}{108}\times(-1)\approx2.9907.
\]
其实, 其精确值为 $\sqrt[4]{80}=2.990697\cdots$。

\noindent\textbf{例 4.3.2}\quad 为了使计算的正方形的面积的相对误差不超过 $1\%$, 度量边长所允许的最大相对误差为多少?

\noindent\textbf{解}\quad 设正方形的边长为 $a$, 面积为 $S$, 则有
\[
  S=a^2,\quad |\Delta S|\approx |dS|=|2a\Delta a|,
\]
\[
  \left|\frac{\Delta S}{S}\right|
  \approx
  \left|\frac{2a\Delta a}{a^2}\right|
  =
  2\left|\frac{\Delta a}{a}\right|.
\]
于是
\[
  \left|\frac{\Delta a}{a}\right|
  \approx
  \frac{1}{2}\left|\frac{\Delta S}{S}\right|
  =
  0.5\%.
\]
这表明度量边长所允许的最大相对误差为 $0.5\%$。

\subsection{一阶微分的形式不变性}

设函数 $y=f(u),u=g(x)$. 根据复合函数的求导法则, 可得复合函数 $y=f(g(x))$ 的微分公式为
\[
  dy=(f(g(x)))'dx=f'(g(x))g'(x)dx.
\]
由于 $du=g'(x)dx$, 代入上式就得到了它的等价表示形式
\[
  dy=f'(u)du.
\]
这里 $u=g(x)$ 是 $x$ 的函数. 但我们发现, 它与 $u$ 为自变量的函数 $f(u)$ 的微分形式
\[
  dy=f'(u)du
\]
一模一样. 换句话说, 对 $f(u)$ 进行微分时, 不管 $u$ 是因变量还是自变量, 所得结果具有相同的形式, 这也就是所谓的\textbf{一阶微分的形式不变性}. 当然它们的意义是不一样的, $u$ 是自变量时, $du=\Delta u$, 而当 $u$ 是函数时, 一般来说 $du$ 与 $\Delta u$ 是不相同的。

这是一阶微分的一个非常重要的性质, 有了这个 ``形式不变性'' 作保证, 我们拿到一个函数 $y=f(u)$ 之后, 就可以按 $u$ 是自变量去求它的微分 $f'(u)du$, 而无须顾忌 $u$ 是自变量, 还是因变量. 这使得微分运算在很多场合比求导数运算优越得多。

例如, 设 $y=f(x)$ 在 $x_0$ 可微, 且存在反函数 $x=f^{-1}(y)$, 再假定 $f'(x_0)\ne0$. 我们先对 $y=f(x)$ 两边微分, 得 $dy=f'(x_0)dx$. 然后将 $x$ 看成因变量, $y$ 看成自变量, 从前面的等式便可导出
\[
  \frac{dx}{dy}=\frac{1}{f'(x_0)},
\]
这正是反函数的求导数公式。

再如, 求参数方程
\[
  \begin{cases}
    x=x(t),\\
    y=y(t)
  \end{cases}
\]
所确定的函数 $y=f(x)$ 的导数. 我们可以先对参数方程分别微分得
\[
  \begin{cases}
    dx=x'(t)dt,\\
    dy=y'(t)dt.
  \end{cases}
\]
然后利用这两个等式, 便可得到
\[
  \frac{dy}{dx}
  =
  \frac{y'(t)dt}{x'(t)dt}
  =
  \frac{y'(t)}{x'(t)},
\]
这正是参数式函数的求导公式。

利用前面得到的基本初等函数的导数公式, 立即可得它们的微分公式为:
\[
\begin{array}{ll}
  (1)\ d(C)=0dx; & (2)\ d(x^\alpha)=\alpha x^{\alpha-1}dx;\\[4pt]
  (3)\ d(a^x)=a^x\ln a\,dx\ (a>0); &
  (4)\ d(\ln|x|)=\dfrac{dx}{x};\\[8pt]
  (5)\ d(\sin x)=\cos x\,dx; &
  (6)\ d(\cos x)=-\sin x\,dx;\\[4pt]
  (7)\ d(\tan x)=\dfrac{dx}{\cos^2 x}; &
  (8)\ d(\arcsin x)=\dfrac{dx}{\sqrt{1-x^2}};\\[8pt]
  (9)\ d(\arctan x)=\dfrac{dx}{1+x^2}. &
\end{array}
\]
同样, 由导数的四则运算法则, 可得微分运算的法则. 设 $u,v$ 为 $x$ 的可微函数, 则
\[
\begin{aligned}
  (1)\quad &d(u\pm v)=du\pm dv;\\
  (2)\quad &d(u\cdot v)=u\,dv+v\,du;\\
  (3)\quad &d\left(\frac{u}{v}\right)=\frac{v\,du-u\,dv}{v^2}.
\end{aligned}
\]
利用这些公式和运算法则, 我们就可以求任何初等函数的微分了。

\noindent\textbf{例 4.3.3}\quad 设 $y=\arctan e^x$, 求 $dy$。

\noindent\textbf{解}\quad
\[
  dy
  =
  \frac{1}{1+e^{2x}}de^x
  =
  \frac{e^x}{1+e^{2x}}dx
  =
  \frac{dx}{e^x+e^{-x}}.
\]

\noindent\textbf{例 4.3.4}\quad 设 $y=\dfrac{x-a}{1-ax}\ (|a|<1)$, 求证: 当 $|x|<1$ 时, 有
\[
  \frac{dy}{1-y^2}=\frac{dx}{1-x^2}.
\]

\noindent\textbf{证明}\quad 由于
\[
  dy=\frac{1-a^2}{(1-ax)^2}dx,
\]
\[
  1-y^2
  =
  1-\left(\frac{x-a}{1-ax}\right)^2
  =
  \frac{(1-a^2)(1-x^2)}{(1-ax)^2},
\]
故有
\[
  \frac{dy}{1-y^2}=\frac{dx}{1-x^2}.
\]
此例的结果在双曲几何中起着重要的作用。

\noindent\textbf{例 4.3.5}\quad 求由方程 $\sin y^2=\cos\sqrt{x}$ 所确定的隐函数 $y=y(x)$ 的导数。

\noindent\textbf{解}\quad 对方程
\[
  \sin y^2=\cos\sqrt{x}
\]
两边微分, 得
\[
  2y\cos y^2\,dy
  =
  -\frac{\sin\sqrt{x}}{2\sqrt{x}}\,dx,
\]
由此即得
\[
  \frac{dy}{dx}
  =
  -\frac{\sin\sqrt{x}}{4y\sqrt{x}\cos y^2}.
\]

\noindent\textbf{注}\quad 读者可以对此例中的隐函数的显式表达式
\[
  y=\pm\sqrt{\arcsin(\cos\sqrt{x})}
\]
直接求导数, 并且将运算过程加以比较。

\section{高阶导数与高阶微分}

\subsection{高阶导数}

当知道物体沿直线运动的路程 $s=s(t)$ 需要求加速度时, 我们就会遇到求速度函数的导数, 即求 $s(t)$ 的二阶导数的问题。

若一个函数 $y=f(x)$ 的一阶导数 $f'(x)$ 仍是可导函数, 则我们可求 $(f'(x))'$, 记其为 $f''(x)$ 或 $\dfrac{d^2y}{dx^2}$, 并称之为 $f$ 的\textbf{二阶导数}. 类似地, 可定义\textbf{三阶导数}为 $f'''(x)=(f''(x))'$. 一般地, 当 $n\ge4$ 时, $f(x)$ 的 $n$ 阶导数定义为 $f(x)$ 的 $n-1$ 阶导数的导数, 并记为 $f^{(n)}(x),y^{(n)}$, 或 $\dfrac{d^ny}{dx^n}$。

对于一些较简单的函数, 我们可以导出它们的高阶导数的公式. 基本方法是, 先求出前几阶导数, 发现规律, 总结出一般的公式, 然后再加以证明。

\noindent\textbf{例 4.4.1}\quad 设 $y=e^{ax}$, 求 $y^{(n)}$。

\noindent\textbf{解}\quad 由于
\[
  y'=ae^{ax},\quad y''=a^2e^{ax},\quad y'''=a^3e^{ax},
\]
因此, 由高阶导数的定义, 容易用数学归纳法证明
\[
  y^{(n)}=a^ne^{ax}\quad(n=1,2,\cdots).
\]

\noindent\textbf{例 4.4.2}\quad 设 $y=x^\alpha$, 求 $y^{(n)}$。

\noindent\textbf{解}\quad 由于
\[
  y'=\alpha x^{\alpha-1},\quad y''=\alpha(\alpha-1)x^{\alpha-2},
\]
因此容易由数学归纳法证明
\[
  y^{(n)}
  =
  \alpha(\alpha-1)\cdots(\alpha-n+1)x^{\alpha-n}
  \quad(n\ge1).
\]

\noindent\textbf{注}\quad 当 $\alpha$ 为正整数时, 若 $\alpha=n$, 则 $y^{(n)}=n!$; 若 $\alpha<n$, 则 $y^{(n)}=0$. 由此立即导出: 如果 $P_m(x)$ 是 $m$ 次多项式, 且 $n>m$, 则
\[
  P_m^{(n)}(x)=0.
\]

\noindent\textbf{例 4.4.3}\quad 设 $y=\ln(1+x)$, 求 $y^{(n)}$。

\noindent\textbf{解}\quad 由
\[
  y'=\frac{1}{1+x},\quad
  y''=\frac{-1}{(1+x)^2},\quad
  y'''=\frac{2!}{(1+x)^3},\quad
  y^{(4)}=\frac{-3!}{(1+x)^4},
\]
可以看出
\[
  y^{(n)}
  =
  (-1)^{n-1}\frac{(n-1)!}{(1+x)^n}
  \quad(n\ge1),
\]
其中规定 $0!=1$。

事实上, 设
\[
  y^{(k)}
  =
  (-1)^{k-1}\frac{(k-1)!}{(1+x)^k},
\]
则
\[
  y^{(k+1)}
  =
  (-1)^{k-1}(k-1)!\frac{-k(1+x)^{k-1}}{(1+x)^{2k}}
  =
  \frac{(-1)^kk!}{(1+x)^{k+1}}.
\]
这样, 由数学归纳法原理知, 对一切的正整数 $n$ 所得公式都成立。

\noindent\textbf{例 4.4.4}\quad 设 $y=\sin x$, 求 $y^{(n)}$。

\noindent\textbf{解}\quad 直接计算, 得
\[
  y'=\cos x=\sin\left(x+\frac{\pi}{2}\right),
  \quad
  y''=-\sin x=\sin(x+\pi).
\]
现在假定
\[
  y^{(n-1)}
  =
  \sin\left(x+\frac{(n-1)\pi}{2}\right),
\]
则有
\[
  y^{(n)}
  =
  \left[\sin\left(x+\frac{(n-1)\pi}{2}\right)\right]'
  =
  \cos\left(x+\frac{(n-1)\pi}{2}\right)
  =
  \sin\left(x+\frac{n\pi}{2}\right).
\]
同理可得
\[
  (\cos x)^{(n)}
  =
  \cos\left(x+\frac{n\pi}{2}\right).
\]

\noindent\textbf{注}\quad 上例中的两个公式也可以借助欧拉公式形式地导出. 在欧拉公式
\[
  e^{ix}=\cos x+i\sin x
\]
的两边形式地求导, 可得
\[
  (e^{ix})^{(n)}
  =
  (\cos x)^{(n)}+i(\sin x)^{(n)},
\]
而
\[
  (e^{ix})^{(n)}
  =
  i^ne^{ix}
  =
  e^{i\frac{n\pi}{2}}e^{ix}
  =
  e^{i\left(x+\frac{n\pi}{2}\right)}
  =
  \cos\left(x+\frac{n\pi}{2}\right)
  +
  i\sin\left(x+\frac{n\pi}{2}\right),
\]
所以
\[
  (\cos x)^{(n)}
  =
  \cos\left(x+\frac{n\pi}{2}\right),\quad
  (\sin x)^{(n)}
  =
  \sin\left(x+\frac{n\pi}{2}\right).
\]

\noindent\textbf{例 4.4.5}\quad 设 $y=\arctan x$, 求 $y^{(n)}$。

\noindent\textbf{解}\quad 由于 $x=\tan y$, 直接计算, 可得
\[
  y'=\frac{1}{1+x^2}
  =
  \frac{1}{1+\tan^2y}
  =
  \cos^2y,
\]
\[
  y''
  =
  -2\cos y\sin y\cdot y'
  =
  \cos^2y\sin2\left(y+\frac{\pi}{2}\right),
\]
\[
\begin{aligned}
  y'''
  &=
  2\cos^3y(-\sin y)\sin2\left(y+\frac{\pi}{2}\right)
  +
  2\cos^4y\cos2\left(y+\frac{\pi}{2}\right)  \\
  &=
  2\cos^3y\cos\left(2\left(y+\frac{\pi}{2}\right)+y\right)  \\
  &=
  2\cos^3y\sin3\left(y+\frac{\pi}{2}\right).
\end{aligned}
\]
一般地, 由数学归纳法可证
\[
  y^{(n)}
  =
  (n-1)!\cos^ny\sin n\left(y+\frac{\pi}{2}\right).
\]
此外, 由这一公式容易导出
\[
  \left.y^{(2n-1)}\right|_{x=0}
  =
  (-1)^{n-1}(2n-2)!,
  \quad
  \left.y^{(2n)}\right|_{x=0}=0.
\]

\subsection{莱布尼茨公式}

对于两个函数 $u,v$ 的和、差的高阶导数, 我们有
\[
  (u\pm v)^{(n)}=u^{(n)}\pm v^{(n)};
\]
对常数 $c$ 和函数 $u$ 的积, 有
\[
  (cu)^{(n)}=cu^{(n)};
\]
而对于两个函数 $u,v$ 的积, 我们有以下莱布尼茨公式:

\noindent\textbf{定理 4.4.1（莱布尼茨公式）}\quad 若函数 $u$ 和 $v$ 有任意阶导数, 则
\[
  (u\cdot v)^{(n)}
  =
  \sum_{k=0}^{n}C_n^ku^{(n-k)}v^{(k)},
\]
其中
\[
  C_n^k=\frac{n!}{k!(n-k)!},\quad u^{(0)}=u,\quad v^{(0)}=v.
\]

\noindent\textbf{证明}\quad 用数学归纳法证之. 当 $n=1$ 时, 公式显然成立. 假设公式对正整数 $n$ 成立, 我们来证公式对正整数 $n+1$ 也成立. 事实上, 应用归纳法假定, 我们有
\[
\begin{aligned}
  (u\cdot v)^{(n+1)}
  &=
  \left[(u\cdot v)^{(n)}\right]'  \\
  &=
  \sum_{k=0}^{n}C_n^k\left[u^{(n-k)}v^{(k)}\right]'  \\
  &=
  \sum_{k=0}^{n}C_n^k
  \left[u^{(n-k+1)}v^{(k)}+u^{(n-k)}v^{(k+1)}\right]  \\
  &=
  \sum_{k=0}^{n}C_n^ku^{(n-k+1)}v^{(k)}
  +
  \sum_{k=1}^{n+1}C_n^{k-1}u^{(n-k+1)}v^{(k)}  \\
  &=
  C_n^0u^{(n+1)}v^{(0)}
  +
  \sum_{k=1}^{n}(C_n^{k-1}+C_n^k)u^{(n-k+1)}v^{(k)}
  +
  C_n^nu^{(0)}v^{(n+1)}  \\
  &=
  \sum_{k=0}^{n+1}C_{n+1}^ku^{(n+1-k)}v^{(k)},
\end{aligned}
\]
这与牛顿二项式展开公式证明的格式是一致的, 其中最后一步用到了恒等式
\[
  C_n^k+C_n^{k-1}=C_{n+1}^k.
\]

\noindent\textbf{例 4.4.6}\quad 设 $y=\arcsin x$, 求 $y^{(n)}(0)$。

\noindent\textbf{解}\quad 由 $y'=\dfrac{1}{\sqrt{1-x^2}}$, 可得
\[
  \sqrt{1-x^2}\cdot y'=1,
\]
再两边平方得
\[
  (1-x^2)\cdot y'^2=1. \tag{4.4.1}
\]
若对上式用莱布尼茨公式, 则 $y'^2$ 的高阶导数仍不好算. 所以对上式再求导数一次, 得
\[
  (1-x^2)\cdot 2y'y''-2xy'^2=0.
\]
由 (4.4.1) 式看出, 当 $|x|<1$ 时, $y'\ne0$, 化简得
\[
  (1-x^2)y''-xy'=0.
\]
对上式两边求 $(n-2)$ 阶导数, 并应用莱布尼茨公式, 得
\[
\begin{aligned}
  &(1-x^2)y^{(n)}
  +(n-2)(-2x)y^{(n-1)}
  +\frac{(n-2)(n-3)}{2}\cdot(-2)y^{(n-2)}  \\
  &\qquad
  -xy^{(n-1)}-(n-2)y^{(n-2)}=0.
\end{aligned}
\]
将 $x=0$ 代入上式, 就有
\[
  y^{(n)}(0)
  -(n-2)(n-3)y^{(n-2)}(0)
  -(n-2)y^{(n-2)}(0)=0,
\]
由此即得递推公式
\[
  y^{(n)}(0)=(n-2)^2y^{(n-2)}(0).
\]
又因 $y^{(0)}(0)=0$, $y^{(1)}(0)=1$, 所以我们有 $y^{(2n)}(0)=0$, 且
\[
\begin{aligned}
  y^{(2n+1)}(0)
  &=
  (2n-1)^2y^{(2n-1)}(0)  \\
  &=
  (2n-1)^2(2n-3)^2y^{(2n-3)}(0)  \\
  &=\cdots  \\
  &=
  (2n-1)^2(2n-3)^2\cdots3^2\cdot1^2\cdot y^{(1)}(0)  \\
  &=
  [(2n-1)!!]^2.
\end{aligned}
\]

\noindent\textbf{例 4.4.7}\quad 设 $y=\tan x$, 求 $y^{(n)}(0)$, $n=1,2,\cdots,8$。

\noindent\textbf{解}\quad 因奇函数的导数为偶函数, 偶函数的导数为奇函数（参考习题四第 4 题）, 以及奇函数在原点的值为零, 所以有
\[
  y(0)=y^{(2)}(0)=y^{(4)}(0)=y^{(6)}(0)=y^{(8)}(0)=0.
\]
为了求奇数阶导数在原点的值, 我们设法建立递推公式. 等式
\[
  y\cdot\cos x=\sin x
\]
两边求 $(2n+1)$ 阶导数, 并利用莱布尼茨公式, 得
\[
  \sum_{k=0}^{2n+1}
  C_{2n+1}^ky^{(2n+1-k)}
  \cos\left(x+\frac{k\pi}{2}\right)
  =
  \sin\left(x+\frac{2n+1}{2}\pi\right),
\]
将 $x=0$ 代入上式, 即得递推公式
\[
  \sum_{k=0}^{2n+1}
  C_{2n+1}^ky^{(2n+1-k)}(0)\cos\frac{k\pi}{2}
  =
  (-1)^n,
\]
即
\[
\begin{aligned}
  &y^{(2n+1)}(0)
  -C_{2n+1}^{2}y^{(2n-1)}(0)
  +C_{2n+1}^{4}y^{(2n-3)}(0)+\cdots \\
  &\quad
  +(-1)^{n-1}C_{2n+1}^{2n-2}y^{(3)}(0)
  +(-1)^nC_{2n+1}^{2n}y'(0)
  =
  (-1)^n.
\end{aligned}
\]
将 $n=1$ 及 $y'(0)=1$ 代入上式, 得
\[
  y^{(3)}(0)-\frac{3\cdot2}{2!}=-1,
\]
解出 $y^{(3)}(0)=2$。

在递推公式中令 $n=2$, 得
\[
  y^{(5)}(0)-\frac{5\cdot4}{2!}\cdot2
  +
  \frac{5\cdot4\cdot3\cdot2}{4!}\cdot1
  =
  1.
\]
解出 $y^{(5)}(0)=16$。

再在递推公式中令 $n=3$, 得
\[
  y^{(7)}(0)-\frac{7\cdot6}{2!}\cdot16
  +
  \frac{7\cdot6\cdot5\cdot4}{4!}\cdot2
  -
  \frac{7!}{6!}\cdot1
  =
  -1.
\]
解出 $y^{(7)}(0)=272$。

附带有
\[
  \frac{y'(0)}{1!}=1,\quad
  \frac{y^{(3)}(0)}{3!}=\frac13,\quad
  \frac{y^{(5)}(0)}{5!}=\frac{2}{15},\quad
  \frac{y^{(7)}(0)}{7!}=\frac{17}{315}.
\]

\subsection{一般函数的高阶导数}

对于一般没有什么特性的函数, 通常很难归纳出其 $n$ 阶导数的公式, 只能逐次地求导. 特别对于复合函数、反函数、隐函数和参数式函数求高阶导数更是如此。

\noindent\textbf{例 4.4.8}\quad 设函数 $y=y(x)$ 是由方程 $\dfrac{x^2}{a^2}+\dfrac{y^2}{b^2}=1\ (y>0)$ 给出, 试求 $y''$。

\noindent\textbf{解法 1}\quad 先从方程解出
\[
  y=\frac{b}{a}\sqrt{a^2-x^2},
\]
然后直接计算, 可得
\[
  y'
  =
  -\frac{b}{a}(a^2-x^2)^{-\frac12}x
  =
  -\frac{bx}{a\sqrt{a^2-x^2}},
\]
\[
  y''
  =
  -\frac{b}{a}(a^2-x^2)^{-\frac32}x^2
  -
  \frac{b}{a}(a^2-x^2)^{-\frac12}
  =
  -\frac{b^4}{a^2y^3}.
\]

\noindent\textbf{解法 2}\quad 先将方程表示为参数式
\[
  \begin{cases}
    x=a\cos t,\\
    y=b\sin t,
  \end{cases}
\]
再直接计算, 有
\[
  y'_x=-\frac{b\cos t}{a\sin t}=-\frac{b}{a}\cot t,
\]
\[
\begin{aligned}
  y''_{xx}
  &=
  -\frac{d(y'_x)}{dx}
  =
  -\frac{b}{a}\frac{(\cot t)'}{x'_t}  \\
  &=
  \frac{b}{a\sin^2t}\cdot\frac{-1}{a\sin t}
  =
  \frac{-b}{a^2\sin^3t}
  =
  -\frac{b^4}{a^2y^3}.
\end{aligned}
\]

\noindent\textbf{解法 3}\quad 直接对方程两边关于 $x$ 求导数, 得
\[
  \frac{2x}{a^2}+\frac{2yy'}{b^2}=0. \tag{4.4.2}
\]
由此解得
\[
  y'=-\frac{b^2x}{a^2y}.
\]
再对方程 (4.4.2) 两边关于 $x$ 求导数, 得
\[
  \frac1{a^2}+\frac{y'^2+yy''}{b^2}=0.
\]
由此解出 $y''$ 并且将 $y'$ 的表达式代入, 得
\[
  y''=-\frac{b^4}{a^2y^3}.
\]
比较上面的三种解法, 不难发现, 似乎隐函数求导法要简单一些。

\subsection{高阶微分}

若函数 $f(x)$ 在区间 $(a,b)$ 内可微, 则有 $\mathrm df=f'(x)\mathrm dx$, 其中 $f'(x)$ 是 $x$ 的函数, 而 $\mathrm dx$ 是与 $x$ 无关的一个量. 这样, 把一阶微分 $\mathrm df$ 看成 $x$ 的函数, 再求一次微分 $\mathrm d(\mathrm df)$, 就称之为 $f(x)$ 的二阶微分, 记为 $\mathrm d^2f$, 即
\[
  \mathrm d^2f
  =
  \mathrm d(\mathrm df)
  =
  (f'(x)\mathrm dx)'\mathrm dx
  =
  f''(x)\mathrm dx\mathrm dx
  =
  f''(x)\mathrm dx^2,
\]
其中 $\mathrm dx^2=\mathrm dx\mathrm dx$（读者应注意 $\mathrm dx^2\ne \mathrm d(x^2)=2x\mathrm dx$）。

一般地, 我们定义 $n$ 阶微分为
\[
\begin{aligned}
  \mathrm d^nf
  &=
  \mathrm d(\mathrm d^{n-1}f)
  =
  (f^{(n-1)}(x)\mathrm dx^{n-1})'\mathrm dx  \\
  &=
  f^{(n)}(x)\mathrm dx^{n-1}\mathrm dx
  =
  f^{(n)}(x)\mathrm dx^n.
\end{aligned}
\]
由此即可知道, 为什么将 $f^{(n)}(x)$ 记为 $\dfrac{\mathrm d^nf}{\mathrm dx^n}$ 的缘故了。

由定义可知, 求高阶微分本质上就是求高阶导数, 因此我们不再详细讨论, 但值得指出的是高阶微分是不具有形式不变性的。

事实上, 设 $y=f(u)$, 当 $u$ 为自变量时, 我们有 $\mathrm d^2y=f''(u)\mathrm du^2$; 但当 $u=g(x)$ 是 $x$ 的函数时, 则
\[
  \mathrm df=f'(u)g'(x)\mathrm dx,
\]
\[
  \mathrm d^2f
  =
  f''(u)(g'(x))^2\mathrm dx^2+f'(u)g''(x)\mathrm dx^2
  =
  f''(u)\mathrm du^2+f'(u)\mathrm d^2u,
\]
这比 $u$ 是自变量时多了一项 $f'(u)\mathrm d^2u$! 因此, 当 $\mathrm d^2u\ne0$, 即 $g''(x)\ne0$ 时, 二阶微分就不具有形式不变性了。

\section*{习题四}

\noindent 1. 讨论函数 $f(x)$ 和 $g(x)$ 在 $x=0$ 处的可导性, 其中
\[
  f(x)=
  \begin{cases}
    -x, & x\in\mathbb R\setminus\mathbb Q,\\
    x, & x\in\mathbb Q;
  \end{cases}
\]
\[
  g(x)=
  \begin{cases}
    -x^2, & x\in\mathbb R\setminus\mathbb Q,\\
    x^2, & x\in\mathbb Q.
  \end{cases}
\]

\noindent 2. 设函数 $f(x)$ 在 $(a,b)$ 上有定义, 且在点 $x_0\in(a,b)$ 处可导, 并假定序列 $\{x_n\}$ 和 $\{y_n\}$ 满足
\[
  a<x_n<x_0<y_n<b,\quad n=1,2,\cdots,
\]
且有
\[
  \lim_{n\to\infty}x_n=x_0=\lim_{n\to\infty}y_n.
\]
证明:
\[
  \lim_{n\to\infty}\frac{f(y_n)-f(x_n)}{y_n-x_n}=f'(x_0).
\]

\noindent 3. 设函数 $f(x)$ 在 $(-1,1)$ 上可导, 且满足 $|f(x)|\le|\sin x|$, 求证 $|f'(0)|\le1$。

\noindent 4. 求证:
\begin{enumerate}
  \item 可导的偶函数, 其导数为奇函数;
  \item 可导的奇函数, 其导数为偶函数;
  \item 可导的周期函数, 其导数为相同周期的周期函数。
\end{enumerate}

\noindent 5. 讨论下列函数在指定点 $x_0$ 处的可导性:
\begin{enumerate}
  \item $y=|(x-1)(x-2)^2(x-3)^3|,\quad x_0=1,2,3$;
  \item
  $
  y=
  \begin{cases}
    \dfrac14(x-1)(x+1)^2, & |x|\le1,\\
    |x|-1, & |x|>1,
  \end{cases}
  \quad x_0=-1,0,1
  $。
\end{enumerate}

\noindent 6. 设函数 $f(x)$ 在 $[-1,1]$ 上有定义, $f'(0)$ 存在, 求
\[
  \lim_{n\to\infty}
  \left[
    f\left(\frac1{n^2}\right)
    +
    f\left(\frac2{n^2}\right)
    +\cdots+
    f\left(\frac n{n^2}\right)
    -
    nf(0)
  \right].
\]

\noindent 7. 求下列序列的极限:
\begin{enumerate}
  \item
  $\displaystyle
  \lim_{n\to\infty}
  \left[
    \sin\frac1{n^2}+\sin\frac2{n^2}+\cdots+\sin\frac n{n^2}
  \right]
  $。
  \item
  $\displaystyle
  \lim_{n\to\infty}
  \left(1+\frac1{n^2}\right)
  \left(1+\frac2{n^2}\right)\cdots
  \left(1+\frac n{n^2}\right)
  $。
\end{enumerate}

\noindent 8. 设 $P_n(x)$ 为最高次系数为 1 的 $n(n\ge1)$ 次多项式, $M$ 为 $P_n(x)=0$ 的最大实根, 求证 $P'_n(M)\ge0$。

\noindent 9. 设函数 $\varphi$ 在点 $a$ 连续, $f(x)=|x-a|\varphi(x)$. 求 $f$ 在点 $a$ 的左右导数, 问在什么条件下 $f(x)$ 在点 $a$ 可导?

\noindent 10. 给定曲线 $y=x^2+5x+4$。
\begin{enumerate}
  \item 确定 $b$, 使直线 $y=3x+b$ 为该曲线的切线;
  \item 确定 $m$, 使直线 $y=mx$ 为该曲线的切线。
\end{enumerate}

\noindent 11. 给定曲线 $y=x^3$ 和直线 $y=px-q$, 其中 $p,q$ 为实数, 且 $p>0$。
\begin{enumerate}
  \item $p$ 给定后, 试确定 $q$, 使 $y=px-q$ 是 $y=x^3$ 的切线;
  \item 利用图形求出方程 $x^3-px+q=0$ 有三个不同实根的条件。
\end{enumerate}

\noindent 12。
\begin{enumerate}
  \item 确定 $m$, 使 $y=mx$ 为曲线 $y=\ln x$ 的切线;
  \item 利用图形求出方程 $\ln x-mx=0$ 有两个实根的条件。
\end{enumerate}

\noindent 13. 问抛物线 $y=x^2-2x-1$ 在哪一点的切线垂直于直线 $x+2y-1=0$?

\noindent 14. 求下列函数的导数:
\begin{enumerate}
  \item $y=\sqrt{x}+\dfrac1x-2x^3$;
  \item $y=x^22^x$;
  \item $y=\dfrac{2x-x^2}{1+x+x^2}$;
  \item $y=\dfrac{x\sin x+\cos x}{x\sin x-\cos x}$;
  \item $y=\dfrac1{\sqrt[3]{x}}+\sqrt[3]{x}$;
  \item $y=\dfrac1{1+\sqrt{x}}-\dfrac1{1-\sqrt{x}}$;
  \item $y=x^3\ln x-\dfrac1n x^n$;
  \item $y=\left(x+\dfrac1x\right)\ln x$;
  \item $y=\sec x$;
  \item $y=\dfrac{\cos x}{x^4}\ln\dfrac1x$。
\end{enumerate}

\noindent 15. 确定常数 $a,b$, 使函数
\[
  f(x)=
  \begin{cases}
    ax+b, & x>1,\\
    x^2, & x\le1
  \end{cases}
\]
有连续的导函数。

\noindent 16. 求下列函数的导数:
\begin{enumerate}
  \item $y=x(a^2+x^2)\sqrt{a^2-x^2}$;
  \item $y=\sqrt[3]{\dfrac{1+x^3}{1-x^3}}$;
  \item $y=\ln(\ln x)$;
  \item $y=\dfrac1{2a}\ln\left|\dfrac{a+x}{a-x}\right|$;
  \item $y=\ln(x+\sqrt{a+x^2})$;
  \item $y=\ln\tan\dfrac x2$;
  \item $y=\cos^3x-\cos3x$;
  \item $y=\sin^nx\cos nx$;
  \item $y=\cos(\cos\sqrt{x})$;
  \item $y=\dfrac{\sin^2x}{\sin x^2}$;
  \item $y=(e^x+e^{-x})^2$;
  \item $y=\ln\sqrt{\dfrac{1+\cos x}{1-\cos x}}$。
\end{enumerate}

\noindent 17. 设在 $x=1$ 处有
\[
  \frac{d}{dx}f(x^2)=\frac{d}{dx}f^2(x),
\]
求证 $f'(1)=0$ 或 $f(1)=1$。

\noindent 18. 给定可导函数 $f(x)>0$, 证明函数 $y=f(x)\sin ax\ (a>0)$ 永远夹在两条振幅曲线 $y=\pm f(x)$ 之间, 并与之相切。

\noindent 19. 求下列函数的导数:
\begin{enumerate}
  \item $y=x^x\ (x>0)$;
  \item $y=x^{\tan x}\ (x>0)$;
  \item $y=x^{\ln x}\ (x>0)$;
  \item $y=(1+x)^{\frac1x}\ (x>0)$;
  \item $y=x^{x^x}\ (x>0)$;
  \item $y=e^{-\frac{\sin x}{x^2}}\ (x\ne0)$。
\end{enumerate}

\noindent 20. 求下列函数的导数:
\begin{enumerate}
  \item $y=\arcsin\sqrt{1-x^2}$;
  \item $y=\arccos(\sin x)$;
  \item $y=x\arctan x-\dfrac12\ln(1+x^2)$;
  \item $y=\arctan\dfrac{\sqrt{1-x^2}-1}{x}+\arctan\dfrac{2x}{1-x^2}$;
  \item $y=\arctan(\tan^2x)$;
  \item $y=\left(\dfrac ab\right)^x\left(\dfrac bx\right)^a\left(\dfrac xa\right)^b\ (a,b>0)$;
  \item $y=\ln\left(\arccos\dfrac1{\sqrt{x}}\right)$;
  \item $y=\ln(e^x+\sqrt{1+e^{2x}})$;
  \item $y=\dfrac x2\sqrt{a^2-x^2}+\dfrac{a^2}{2}\arctan\dfrac xa\ (a>0)$;
  \item $y=\dfrac1{4\sqrt2}\ln\dfrac{x^2+\sqrt2x+1}{x^2-\sqrt2x+1}-\dfrac1{2\sqrt2}\arctan\dfrac{\sqrt2x}{x^2-1}$。
\end{enumerate}

\noindent 21. 求证:
\[
  \lim_{x\to a}\frac{f(x)-b}{x-a}=A
\]
的充分必要条件是
\[
  \lim_{x\to a}\frac{e^{f(x)}-e^b}{x-a}=e^bA.
\]

\noindent 22. 设曲线的参数表示: $x=a\cos^3t,\ y=a\sin^3t\ (a>0)$。
\begin{enumerate}
  \item 求 $y'(x)$;
  \item 证明曲线的切线被坐标轴所截的长度为一常数。
\end{enumerate}

\noindent 23. 求椭圆 $\dfrac{x^2}{a^2}+\dfrac{y^2}{b^2}=1$ 在第一象限中的切线, 使它平行于过 $(0,b)$ 和 $(a,0)$ 的直线。

\noindent 24. 证明曲线
\[
  \begin{cases}
    x=a(\cos t+t\sin t),\\
    y=a(\sin t-t\cos t)
  \end{cases}
  \quad(a>0)
\]

\end{document}
